\chapter{Enunciado del alcance}

\begin{quotation}[Psicólogo clínico] {Fitzhugh Dodson (1923--1993)}
Sin metas y planes para llegar a ellos, eres como un barco que ha zarpado sin destino.
\end{quotation}


\begin{abstract}
Para gestionar exitosamente un proyecto y cumplir con las metas planteadas al comienzo del mismo, es prioritario establecer y definir el alcance del mismo. 
\end{abstract}

\clearpage

\section{Acta fundacional}

Reunidas en Sevilla el 12 de noviembre de 2024 a las 17:00 que a continuación se detallan.

\begin{enumerate}
\item Francisco José Anillo Carrasco, de nacionalidad española, con DNI 75967897-R, con domicilio en 11300, La Línea de la Concepción, Cádiz, Calle Gabriel Miró, 52.
    \item Daniel Dorante Sánchez, de nacionalidad española, con DNI 24523601-N, con domicilio en 41710, Utrera, Sevilla, Bda. los Salesianos, Bloque 3 Bajo C
\end{enumerate}
    
\textbf{ACUERDAN}
1º Constituir una asociación al amparo de la Ley Orgánica 3/2009, del 3 de abril, sobre modificaciones estructurales de las sociedades mercantiles, que se denominará  \textbf(FADD).
2º Aprobar los estatutos que se incorporarán en este Acta fundacional como anexo por los cuales se va a regir la entidad, que fueron leídos en este mismo acto y aprobados por la unanimidad de los reunidos.
3º Designar a la Junta Directiva de la presente entidad, cuya composición será la siguiente:
\begin{itemize}
    \item Coordinación del Proyecto: Francisco José Anillo Carrasco.
    \item Subcoordinación del Proyecto: Daniel Dorante Sánchez.
\end{itemize}
4º Consentir a la Administración encargada de la inscripción registral para que sean comprobados los datos de la entidad firmantes (RD 522/2006, del 8 de abril)

Y sin más asuntos que tratar se levanta la sesión siendo las 18:30 horas del 12 de noviebre de 2024.

\clearpage
\section{Anexo I Estatutos de la sociedad}

TÍTULO I. DISPOSICIONES GENERALES\\
\textbf{Artículo 1.} Denominación y naturales

Se constituye en la ciudad de Sevilla, una sociedad mercantil con denominación \textbf{(FA-DD)}, con personalidad jurídica y plena capacidad de obrar y que con carácter general, se regirá por la legislación vigente reguladora de régimen jurídico de las sociedades y por los vigentes Estatutos.

\textbf{Articulo 2.} Domicilio y ámbito de actuación:

El domicilio de la Sociedad se establece en: 
Escuela Técnica Superior de Ingeniería Informática. Universidad de Sevilla. Avenida Reina Mercedes S/N. C.P. 41012 - Sevilla, Sevilla.

\textbf{Artículo 3.} Duración
La sociedad se constituye por un periodo de 12 meses tomando como fecha de vencimiento el miércoles 12 de noviembre de 2025.
\textbf{Artículo 4.} Fines y actividades
Son fines de la sociedad:
\begin{enumerate}
    \item El desarrollo de un servicio de asignación y optimización de interecambios de grupos de la Escuela Técnico Superior de Ingeniería Informática, que sea ejecutable en los navegadores habituales.
    \item El desarrollo e implantación de mecanismos y herramientas web que le den soporte y continuidad al servicio previamente definido.
    \item La promoción de la práctica del trabajo en grupo para potenciar esta habilidad transversal en los futuros ingenieros informáticos egresados de la Universidad de Sevilla.
    \item El fomento del trabajo sobre nuevas tecnologías vigentes o con proyección en el mercado laboral meidante el aprendizaje autónomo de estas.
\end{enumerate}

TÍTULO II. DE LOS SOCIOS\\

\textbf{Artículo 5.} Del socio
Serán socios de la Sociedad aquellas personas aprobadas por la Junta Directiva de la Sociedad e inscritos en el registro de socios del "Documento del Proyecto".
\textbf{Artículo 6.} De los derechos del socio

Son derechos del socio los siguientes:

\begin{itemize}
    \item La asistencia a las reuniones realizadas por la Sociedad.
    \item La participación en las decisiones tomadas por la Sociedad.
    \item La libre emisión del voto en las distintas votaciones realizadas por la Sociedad.
    \item La solicitud de un periodo de baja carga de trabajo por motivos justificados con la debida antelación oportuna.
    \item La solicitud de ayuda al resto de integrantes de la sociedad ante una dificultad en el desarrollo del proyecto.
    \item La solicitud de cambio de tareas asignadas con otro integrante de la sociedad, previo acuerdo y aprobación de las partes, siempre que no se haya comenzado con el desarrollo de ninguna de las partes.
    
\end{itemize}

\textbf{Artículo 7.} De las obligaciones del socio

Son las obligaciones del socio las siguientes:

\begin{itemize}
    \item La asistencia a las reuniones realizadas por la sociedad.
    \item La realización de las tareas asignadas por la Jefatura del Proyecto.
    \item La correcta comunicación con el resto de los integrantes promoviendo el buen trato así como el satisfactorio desarrollo del proyecto.
    \item La comunicación con la Sociedad de toda aquella situación o acción que pueda suponer un perjuicio para el proyecto.
\end{itemize}

\textbf{Artículo 8.} De las sanciones del socio

Las sanciones al Socio podrán comprender desde amonestaciones verbales o escritas, hasta la expulsión de la Sociedad. Serán propuestas por la Junta directiva de la Sociedad, bajo demanda de un integrante de esta.
Será responsabilidad de la Sociedad la aprobación, por mayoría simple, de las medidas propuestas por la Junta Directiva oídas las partes implicadas y siempre atendiendo a la situación actual y a los posibles motivos o avisos que haya dado el integrante señalado para el retraso o el no cumplimiento de alguno de los puntos.

TÍTULO III. DE LA COORDINACIÓN DE LA SOCIEDAD\\
\textbf{Artículo 9.} De la Junta Directiva
La Junta Directiva estará compuesta por dos personas: la Coordinación del proyecto y la Jefatura del proyecto. 

Serán funciones de la Junta Directiva:
\begin{itemize}
    \item Asegurar la buena dirección y el éxito de la Sociedad. 
    \item La toma de decisiones consideradas por esta como críticas para la Sociedad. 
    \item Los desempates de mutuo acuerdo cuando la Sociedad se encuentre paralizada por una diferencia equivalente de opiniones de sus socios.
\end{itemize}

\textbf{Artículo 10.} De la Coordinación del proyecto
La Coordinación del proyecto será ostentada por una persona designada por el conjunto de la Sociedad.

Corresponden a la Coordinación del proyecto las siguientes funciones: 
\begin{itemize}
    \item La coordinación de la Sociedad con cuantas sociedades se considere necesario.
    \item La transmisión de toda la información recibida por las otras coordinaciones de proyecto a la Sociedad.
    \item La defensa de los derechos de la Sociedad frente a otras sociedades.
\end{itemize}

\textbf{Artículo 10.} De la Subcoordinación del proyecto
La Coordinación del proyecto será ostentada por una persona designada por el conjunto de la Sociedad.

Corresponden a la Subcoordinación del proyecto las siguientes funciones: 
\begin{itemize}
    \item Las funciones que la Junta Directiva estime necesarias.
    \item La toma de las responsabilidades de la Coordinación del proyecto en caso de indisposición temporal o perpetua de la Coordinación del proyecto.
\end{itemize}

